\anexo{Cronograma de Actividades}

El proyecto se organiza en actividades principales distribuidas a lo largo de las cinco entregas pautadas por la cátedra durante el ciclo lectivo 2026: la propuesta inicial (Entrega~1); los antecedentes —marco teórico y estado del arte— (Entrega~2); la investigación con usuarios y el análisis de negocio (Entrega~3); el diseño y el desarrollo del prototipo (Entrega~4); y la validación experimental junto con las conclusiones finales (Entrega~5). Siguiendo una dinámica de trabajo incremental, el avance se planifica y se revisa antes de cada instancia de evaluación, lo que permite ajustar el alcance en función de la retroalimentación recibida.

La Figura~\ref{fig:cronograma} presenta el cronograma en formato de diagrama de Gantt organizado por entregas, mientras que la Tabla~\ref{tab:plan-actividades} detalla cada actividad junto con la entrega en la que se desarrolla. Las actividades correspondientes a las Entregas~1 y~2 se encuentran completadas a la fecha del presente documento; las restantes reflejan la planificación prevista para el resto del año.

\begin{figure}[ht]
    \centering
    \resizebox{\textwidth}{!}{%
    \begin{ganttchart}[
        hgrid,
        vgrid,
        x unit=2.4cm,
        y unit title=0.7cm,
        y unit chart=0.6cm,
        title height=1,
        bar height=0.55,
        bar/.append style={fill=blue!40},
        bar label font=\normalsize,
        title label font=\normalsize\bfseries
    ]{1}{5}
        \gantttitle{Entregas}{5} \\
        \gantttitle{E1 (Propuesta)}{1}
        \gantttitle{E2 (25\%)}{1}
        \gantttitle{E3 (50\%)}{1}
        \gantttitle{E4 (75\%)}{1}
        \gantttitle{E5 (Final)}{1} \\
        \ganttbar{Definición del proyecto}{1}{1} \\
        \ganttbar{Marco teórico}{2}{2} \\
        \ganttbar{Estado del arte}{2}{2} \\
        \ganttbar{User research}{3}{3} \\
        \ganttbar{Competencia y modelo de negocio}{3}{3} \\
        \ganttbar{Construcción del dataset argentino}{3}{4} \\
        \ganttbar{Arquitectura y tecnologías}{4}{4} \\
        \ganttbar{Desarrollo del prototipo}{4}{5} \\
        \ganttbar{Experimentación y validación}{5}{5} \\
        \ganttbar{Redacción y revisión del documento}{1}{5}
    \end{ganttchart}%
    }
    \caption{Cronograma de actividades del proyecto organizado por entregas (año 2026).}
    \label{fig:cronograma}
\end{figure}

\begin{table}[ht]
    \centering
    \caption{Plan de actividades del proyecto por entrega. Fuente: elaboración propia.}
    \label{tab:plan-actividades}
    \begin{tabular}{c l c}
        \hline
        \textbf{ID} & \textbf{Actividad} & \textbf{Entrega} \\
        \hline
        A-01 & Definición del tema, alcance y objetivos & 1 \\
        A-02 & Marco teórico & 2 \\
        A-03 & Estado del arte & 2 \\
        A-04 & User research (entrevistas, encuestas y \textit{personas}) & 3 \\
        A-05 & Análisis de competencia y modelo de negocio & 3 \\
        A-06 & Arquitectura, datasets y tecnologías & 4 \\
        A-07 & Desarrollo del prototipo (modelos y módulos) & 4 \\
        A-08 & Experimentación, validación y conclusiones & 5 \\
        \hline
    \end{tabular}
\end{table}
